\section{CGA numerical experiments and finite-window diagnostics}
\label{sec:numerics}

\subsection{Protocol, representative problems, and error quantities}

All problems are posed on \(\Omega=(0,1)^d\), with \(d=1\) or \(2\), and use manufactured
profiles. The main suite uses smooth profiles; the final row is explicitly a low-order activation
case \((p,k)=(5,1)\), so its regularity statement is kept separate. The exact profiles, loads, homogeneous natural boundary data, feasible spaces, and
casewise configurations are indexed in Supplement S1 and \texttt{manifest/experiment\_index\_v2.csv}. The exact solution is
used to construct the load and to evaluate errors, but not in the greedy selection or coefficient
optimization. Pure and regularized \(p\)-problems use a zero-mean representative, and their atoms are
centered before normalization.

For a raw atom \(\varphi\), let \(P_\circ\varphi\) denote centering in the pure and regularized cases and
the identity in the other cases. With \(Q_{\rm tr}\) denoting the training quadrature, the experimental
atom is \(g_Q=P_\circ\varphi/\mathcal N_Q(P_\circ\varphi)\), where
\[
 \mathcal N_Q(g)=
 \begin{cases}
 [Q_{\rm tr}(\abs g^2+\abs{\nabla g}^2)]^{1/2},&\text{linear and semilinear cases},\\
 [Q_{\rm tr}(\abs{\nabla g}^p)]^{1/p},&\text{pure and regularized \(p\)-cases},\\
 [Q_{\rm tr}(\abs g^p+\abs{\nabla g}^p)]^{1/p},&\text{reaction \(p\)-case}.
 \end{cases}
\]
This discrete homogeneous normalization is used for atom construction; it is distinct from the
theoretically normalized continuous dictionary in \Cref{sec:setting} and from the nonlinear
\(V\)-distance used to evaluate an approximation. Consequently, a discrete coefficient norm or
training-quadrature normalization is not by itself an entropy or source bound in the analytical norm.

The coefficient solver is monitored in the stable coordinates supplied by the active feature basis.
If an implementation reports a Euclidean coefficient gradient instead, the active-space Gram matrix
is required to convert it into the dual norm on the active span; if the gradient is evaluated with
the training quadrature, a derivative quadrature term must also be included. A dimensionless
optimizer residual is therefore not by itself a value of either \(\eta_m^{\rm stat}\) or
\(\varepsilon_m^{\rm opt}\) in \Cref{thm:inexact}. The run records give the projected residual, the
common target, the optimizer status, and the energy decrease separately.

At outer attempt \(a\), the available atom with largest derivative score is tested, and all active outer
coefficients are reoptimized if the innovation test succeeds. An accepted step \(m\) is one that passes
the rank, finite-value, stationarity, energy-decrease, and feasibility tests. Its active rank is
\(N_m\), while \(M\) denotes the last accepted step reported for a run; rejected attempts increase
\(a\) but not \(m\) or \(N_m\). In the present runs accepted atoms are linearly independent, so the
plotting abscissa \(N=N_m=m\). The maximum attempted count specified by the configuration and all stopping reasons are
recorded in the configuration and run summaries. The eight cases in \Cref{tab:cases} include a
quadratic calibration, two semilinear energies, three \(p=4\) diffusion structures, a
one-dimensional pure \(p=4\) problem, and a low-regularity \(p=5\) case.

\begin{table}[!htbp]
\centering
\small
\caption{Representative CGA experiments and terminal accepted widths.}
\label{tab:cases}
\begin{tabular}{@{}lcccr@{}}
\toprule
Energy & \(d\) & \(p\) & \(k\) & Accepted/target \\
\midrule
linear reaction--diffusion, multifrequency & 1 & 2 & 3 & 256/256 \\
cubic semilinear, multifrequency & 1 & -- & 3 & 256/256 \\
hyperbolic-sine semilinear & 2 & -- & 3 & 512/512 \\
pure \(p\)-Laplacian, low-frequency profile & 1 & 4 & 3 & 141/256 \\
pure \(p\)-Laplacian & 2 & 4 & 3 & 512/512 \\
regularized \(p\)-Laplacian, \(\varepsilon=0.1\) & 2 & 4 & 3 & 512/512 \\
reaction \(p\)-Laplacian & 2 & 4 & 3 & 512/512 \\
pure \(p\)-Laplacian & 1 & 5 & 1 & 256/256 \\
\bottomrule
\end{tabular}
\end{table}

For the main eight-case suite, the one-dimensional target width is 256 and the
two-dimensional target width is 512. The separate quartic diagnostic runs use
\(N=8\)--32 for calibration and continue through \(N=64\). The
one-dimensional pure \(p=4\) case reaches 141 accepted atoms within 768 attempted selections and
therefore has \(N=128\) as its last dyadic
state. The CGA curves show the accepted states over the reported range. Markers identify the
powers of two used for local-order tables and are the only states joined by the display curves; the
complete accepted-state record is retained in the machine-readable supplement.

The quartic finite-trajectory calculation uses two distinct windows. Constants and source
quantities are calibrated only on the 25 accepted states $N=8,\ldots,32$; the subsequent
states $N=33,\ldots,64$ are a continuation of the same saved prefixes and are not used to
re-estimate those constants. Consequently, a pass count on $32<N\leq64$ is a diagnostic of
whether the calibration remains usable immediately beyond its fitting window, whereas the
$8\leq N\leq32$ quantities are the only ones entering the stated finite-window envelope. The two
sets of transitions are reported separately below.

For the linear and semilinear cases we report the energy gap and
\[
 e_{H^1}(\widehat G_N)=\frac{\norm{\widehat G_N-u^*}_{H^1}}{\norm{u^*}_{H^1}}.
\]
For a \(p\)-model the relative Sobolev error and relative natural distance are
\begin{align*}
 e_{W^{1,p}}(\widehat G_N)&=\frac{\|\widehat G_N-u^*\|_{W^{1,p}}}{\|u^*\|_{W^{1,p}}},\\
 e_V(\widehat G_N)&=\frac{d_{V,\bullet}(\widehat G_N,u^*)}{d_{V,\bullet}(0,u^*)}.
\end{align*}
The benchmark solutions make these denominators nonzero. The energy gap is the absolute quantity
\(\widehat\delta_N=\E(\widehat G_N)-\E(u^*)\).
The dyadic local order is
\[
 \operatorname{ord}_N(e)=\log_2\!\left(\frac{e(N/2)}{e(N)}\right).
\]
These three quantities are kept separate in every figure, table, and subsequent comparison.

\subsection{Conditional finite-window guides}

For \(\relu^k\) atoms, the Hilbert-space approximation exponent associated with the smooth-ridge
entropy estimate is
\[
 \beta_H=\frac12+\frac{2(k-1)+1}{2d}.
\]
The dashed segments in the CGA convergence figures show the conditional powers \(N^{-2\beta_H}\) for
the energy gap and \(N^{-\beta_H}\) for either the \(H^1\) error or the natural \(V\)-distance. For a
standard \(p\)-structure energy, the corresponding comparison power for the \(W^{1,p}\) error is
\(N^{-2\beta_H/p}\), as follows from the energy-to-metric conversion in
\Cref{cor:metric-rates}. The exponent is fixed by the stated theory and is not fitted to the curve.
The vertical position is chosen once from the last up to three positive dyadic states by
setting \(\log C=\operatorname{median}_i(\log e_i+s\log N_i)\), where \(e_i>0\)
are those terminal values and \(s\) is the prescribed exponent. This minimizes the
sum of absolute log deviations for the guide \(CN^{-s}\). The segment starts at \(N=8\) in one dimension and at
\(N=16\) in two dimensions and extends to the largest accepted width shown in the same panel. The
display interval is not a claim that the continuous coverage, source, and transfer assumptions have
been verified on that interval.

\begin{table}[!htbp]
\centering
\small
\caption{Exponents used only for the conditional guide segments in the CGA figures.}
\label{tab:window-orders}
\begin{tabular}{@{}ccccc@{}}
\toprule
\((k,d,p)\) & \(\beta_H\) & Energy guide & \(H^1/V\) guide & \(W^{1,p}\) guide \\
\midrule
\((1,1,5)\) & \(1\) & \(2\) & \(1\) & \(2/5\) \\
\((3,1,4)\) & \(3\) & \(6\) & \(3\) & \(3/2\) \\
\((3,2,4)\) & \(7/4\) & \(7/2\) & \(7/4\) & \(7/8\) \\
\((3,1,{--})\) & \(3\) & \(6\) & \(3\) & {--} \\
\((3,2,{--})\) & \(7/4\) & \(7/2\) & \(7/4\) & {--} \\
\bottomrule
\end{tabular}
\end{table}
\FloatBarrier

\subsection{Linear and cubic calibration}

\Cref{fig:cga-linear-cubic} shows the accepted dyadic states for the linear and cubic multifrequency
problems.  Both curves decrease rapidly after about 32 atoms.  Between \(N=128\) and 256, the
linear energy and \(H^1\) orders are 6.083 and 3.042, whereas the cubic orders are 6.896 and 3.448; the
corresponding terminal relative \(H^1\) errors are \(1.247\times10^{-4}\) and
\(1.215\times10^{-4}\).  These are measured local orders on one resolved interval, not estimates of an
asymptotic exponent.

\begin{figure}[!htbp]
\centering
\includegraphics[width=0.96\textwidth]{cga_linear_cubic.pdf}
\caption{CGA values at accepted powers of two for the linear and cubic problems. Each dashed segment
is the corresponding metric-specific conditional guide of
\Cref{tab:window-orders}, restricted to the declared display window.}
\label{fig:cga-linear-cubic}
\end{figure}

\begin{table}[!htbp]
\centering\small
\caption{Dyadic errors for Linear reaction--diffusion, d=1.}
\label{tab:dyadic-01}
\begin{tabular}{@{}rrrrr@{}}
\toprule
$N$ & Energy gap & Order & Relative $H^1$ error & Order \\ \midrule
32 & $4.277\!\times\!10^{-1}$ & -- & $1.323\!\times\!10^{-1}$ & -- \\
64 & $3.983\!\times\!10^{-3}$ & $6.747$ & $1.277\!\times\!10^{-2}$ & $3.373$ \\
128 & $2.577\!\times\!10^{-5}$ & $7.272$ & $1.027\!\times\!10^{-3}$ & $3.636$ \\
256 & $3.801\!\times\!10^{-7}$ & $6.083$ & $1.247\!\times\!10^{-4}$ & $3.042$ \\
\bottomrule
\end{tabular}
\end{table}

\begin{table}[!htbp]
\centering\small
\caption{Dyadic errors for Cubic semilinear problem, d=1.}
\label{tab:dyadic-03}
\begin{tabular}{@{}rrrrr@{}}
\toprule
$N$ & Energy gap & Order & Relative $H^1$ error & Order \\ \midrule
32 & $3.752\!\times\!10^{-1}$ & -- & $1.239\!\times\!10^{-1}$ & -- \\
64 & $2.238\!\times\!10^{-3}$ & $7.389$ & $9.569\!\times\!10^{-3}$ & $3.694$ \\
128 & $4.301\!\times\!10^{-5}$ & $5.702$ & $1.326\!\times\!10^{-3}$ & $2.851$ \\
256 & $3.610\!\times\!10^{-7}$ & $6.896$ & $1.215\!\times\!10^{-4}$ & $3.448$ \\
\bottomrule
\end{tabular}
\end{table}



The two-dimensional sinh case is reported in the supplementary figures.  Its relative \(H^1\) error
decreases from \(4.074\times10^{-3}\) at \(N=256\) to \(1.990\times10^{-3}\) at \(N=512\), while the
local order decreases from 1.793 on \(128\to256\) to 1.034 on \(256\to512\). The final interval
therefore exhibits slower decay than the preceding interval.
\FloatBarrier

\subsection{Pure and regularized \texorpdfstring{\(p=4\)}{p=4} behavior}

\Cref{fig:cga-p4-representative} presents the one- and two-dimensional pure \(p=4\) cases using energy,
relative \(W^{1,4}\) error, and relative \(V\)-distance. The one-dimensional case reaches relative errors
\(1.579\times10^{-5}\) in \(W^{1,4}\) and \(1.680\times10^{-6}\) in the \(V\)-distance at its last
dyadic state \(N=128\). For the two-dimensional pure \(p=4\) problem, the interval \(128\to256\) has energy and \(V\)-orders 3.410 and
1.703.  On the final interval the \(V\)-distance continues to decrease, but the relative \(W^{1,4}\)
error increases from \(2.494\times10^{-3}\) to \(3.601\times10^{-3}\). Thus the terminal
Sobolev-error behavior is nonmonotone even though the \(V\)-distance decreases.

\begin{figure}[!htbp]
\centering
\includegraphics[width=0.98\textwidth]{cga_p4_representative.pdf}
\caption{CGA values at accepted powers of two for the representative one- and two-dimensional pure \(p=4\) problems. The
three columns show energy gap, relative \(W^{1,4}\) error, and relative \(V\)-distance. Dashed segments
use the metric-specific conditional exponents in \Cref{tab:window-orders}.}
\label{fig:cga-p4-representative}
\end{figure}

\begin{table}[!htbp]
\centering\scriptsize
\caption{Dyadic errors for Pure p=4 problem, d=1, k=3.}
\label{tab:dyadic-17}
\begin{tabular}{@{}rrrrrrr@{}}
\toprule
$N$ & Energy gap & Order & Relative $W^{1,4}$ & Order & Relative $V$ & Order \\ \midrule
8 & $2.363\!\times\!10^{-2}$ & -- & $3.567\!\times\!10^{-2}$ & -- & $1.028\!\times\!10^{-2}$ & -- \\
16 & $3.372\!\times\!10^{-4}$ & $6.131$ & $4.460\!\times\!10^{-3}$ & $2.999$ & $1.241\!\times\!10^{-3}$ & $3.051$ \\
32 & $3.180\!\times\!10^{-6}$ & $6.728$ & $2.869\!\times\!10^{-4}$ & $3.959$ & $1.205\!\times\!10^{-4}$ & $3.364$ \\
64 & $1.441\!\times\!10^{-7}$ & $4.464$ & $1.800\!\times\!10^{-4}$ & $0.672$ & $2.565\!\times\!10^{-5}$ & $2.232$ \\
128 & $6.187\!\times\!10^{-10}$ & $7.864$ & $1.579\!\times\!10^{-5}$ & $3.511$ & $1.680\!\times\!10^{-6}$ & $3.932$ \\
\bottomrule
\end{tabular}
\end{table}

\begin{table}[!htbp]
\centering\scriptsize
\caption{Dyadic errors for Pure p=4 problem, d=2, k=3.}
\label{tab:dyadic-08}
\begin{tabular}{@{}rrrrrrr@{}}
\toprule
$N$ & Energy gap & Order & Relative $W^{1,4}$ & Order & Relative $V$ & Order \\ \midrule
16 & $3.070\!\times\!10^{1}$ & -- & $3.624\!\times\!10^{-1}$ & -- & $3.862\!\times\!10^{-1}$ & -- \\
32 & $6.476\!\times\!10^{0}$ & $2.245$ & $1.528\!\times\!10^{-1}$ & $1.247$ & $1.786\!\times\!10^{-1}$ & $1.113$ \\
64 & $1.479\!\times\!10^{-1}$ & $5.452$ & $3.078\!\times\!10^{-2}$ & $2.311$ & $2.701\!\times\!10^{-2}$ & $2.725$ \\
128 & $1.140\!\times\!10^{-2}$ & $3.697$ & $1.017\!\times\!10^{-2}$ & $1.598$ & $7.452\!\times\!10^{-3}$ & $1.858$ \\
256 & $1.073\!\times\!10^{-3}$ & $3.410$ & $2.494\!\times\!10^{-3}$ & $2.028$ & $2.288\!\times\!10^{-3}$ & $1.703$ \\
512 & $4.975\!\times\!10^{-4}$ & $1.109$ & $3.601\!\times\!10^{-3}$ & $-0.530$ & $1.584\!\times\!10^{-3}$ & $0.531$ \\
\bottomrule
\end{tabular}
\end{table}



To examine the transition from uniformly elliptic regularizations to the pure model on one common
finite window, \Cref{fig:cga-epsilon-scan} compares \(\varepsilon=0,10^{-2},10^{-1},1\) for
\(p=4,d=1,k=3\).  At \(N=32\), the four relative \(W^{1,4}\) errors lie between
\(3.209\times10^{-4}\) and \(3.261\times10^{-4}\), while the relative \(V\)-distances lie between
\(9.518\times10^{-5}\) and \(1.198\times10^{-4}\).  The spread of these values indicates weak
dependence on the regularization over this parameter range and window, and it does not establish
coercivity of the unregularized Hessian on the full energy space.

\begin{figure}[!htbp]
\centering
\includegraphics[width=0.98\textwidth]{cga_epsilon_scan.pdf}
\caption{Pure and regularized \(p=4,d=1,k=3\) CGA trajectories for four values of \(\varepsilon\).
Only dyadic accepted states are connected. All runs use the same target width and evaluator.}
\label{fig:cga-epsilon-scan}
\end{figure}

The regularized and reaction \(p=4\) trajectories are shown in
the supplementary figures.  As in the pure two-dimensional case, their final \(W^{1,4}\) errors
increase while their \(V\)-distances continue to decrease.  The same nonmonotone terminal behavior
therefore appears in all three \(p=4\) structures.
\FloatBarrier

\subsection{Finite-trajectory bounds and fractional innovation for the quartic models}
\label{sec:finite-trajectory-numerics}
\label{sec:fractional-innovation-p4}

We evaluate the estimates of \Cref{sec:c5} on the one-dimensional pure and
regularized \(p=4\), \(k=3\) trajectories, with \(u^*(x)=\cos(2\pi x)\) and
\(\varepsilon=0\) or \(0.1\) (cases \texttt{base\_pure} and \texttt{epsilon\_01}
in Supplement S1). These runs use a larger pool and a smaller accepted-state budget than the
eight cases of \Cref{sec:numerics}. The inner product held fixed during the estimates is
\[
 a_*(v,w)=\int_0^1\big(\varepsilon^2+3|(u^*)'|^2\big)v'w'\dd x.
\]
The calibration range contains the 25 retained states \(N=8,\ldots,32\) and its 24 transitions
\(N\to N+1\), \(N=8,\ldots,31\).  The subsequent finite range contains the 32 transitions
\(N\to N+1\), \(N=32,\ldots,63\), generated by extending the same previously computed iterates to \(N=64\).
The two transition sets are kept separate: all source, comparison, and fractional-innovation
constants are estimated on the calibration range only, while the continuation is used only as a
auxiliary numerical test.
At the starting index, the fixed auxiliary dictionary comprises 463 valid candidates and the first 128
valid reference atoms in the order in which they were generated, drawn with a separate seed. They
need not be linearly independent. Symmetry and the frozen normalization follow the theory, and
Supplement S2 specifies the ordering and the filtering. The reference atoms enter the estimates
only. Source coefficients are constructed from the starting-index residual and are not adjusted to fit
later errors.

The full active-space projections use untruncated QR on a common breakpoint-segmented Gauss
rule of order 20, cross-checked at order 16. The active linear systems have condition numbers
between \(7.7\times10^1\) and \(4.1\times10^3\). The maximum defect in
\cref{eq:c5-quartic} is below \(8.7\times10^{-17}\), and the projection-drop identity
has residual below \(1.5\times10^{-12}\). These are floating-point calculations, not interval
enclosures of the exact integrals or projections.

\paragraph{Source bound and Hessian-transfer conditions.}
\Cref{tab:merged-trajectory-conditions} collects the source data, the quartic comparison
constants, and the initial-value bounds of \Cref{thm:c5-source}, evaluated with
\(\underline\theta_m=\theta_m\), the observed selection ratio. In the table,
\(B=B_{\rm entry}\) and \(q=q_{\rm corr}\). The source energy
bounds at \(N=32\) are \(2.23\times10^{-2}\) and \(2.11\times10^{-2}\), compared with
energy gaps \(1.99\times10^{-6}\) and \(3.01\times10^{-6}\). Thus the source remainder
does not make this bound sharp; the accumulated-innovation contribution remains dominant.
For both models, \(\sigma^2/r_8^2\) is approximately \(6.21\times10^{-5}\), whereas
\(\sigma^2/r_{32}^2\) is \(0.751\) for the pure model and \(0.495\) for the regularized model.

\begin{table}[!htbp]
\centering
\small
\caption{Source and quartic comparison quantities on all retained states \(N=8\)--32. The source is fixed at the starting index; \(q\) and \(\rho\) are calibration-range maxima on \(N=8\)--32. Values are floating-point estimates.}
\label{tab:merged-trajectory-conditions}
\begin{tabularx}{\linewidth}{@{}Xrr@{}}
\toprule
Quantity & Pure \(p=4\) & \(\varepsilon=0.1\) \\
\midrule
Source budget \(B\) & \(4.361\times10^{1}\) & \(4.362\times10^{1}\) \\
Source remainder \(\sigma\) & \(1.726\times10^{-3}\) & \(1.728\times10^{-3}\) \\
\(\mathcal W_{24}\) & \(1.912\times10^{4}\) & \(2.913\times10^{4}\) \\
\(q=\max D_N/r_N\) & \(4.733\times10^{-2}\) & \(4.715\times10^{-2}\) \\
\(\rho=\max |R_N|/(r_N^2+D_N^2)\) & \(2.006\times10^{-2}\) & \(1.997\times10^{-2}\) \\
Computed \(r_{32}^2\) & \(3.969\times10^{-6}\) & \(6.025\times10^{-6}\) \\
Initial-value source bound for \(r_{32}^2\) & \(4.280\times10^{-2}\) & \(4.057\times10^{-2}\) \\
Computed energy gap at \(N=32\) & \(1.985\times10^{-6}\) & \(3.013\times10^{-6}\) \\
Initial-value source energy bound & \(2.231\times10^{-2}\) & \(2.114\times10^{-2}\) \\
\bottomrule
\end{tabularx}
\end{table}


For comparison with the finite-range discussion, the minimum active-space generalized Hessian
eigenvalues relative to the \(H^1\) Gram matrix are \(0.1328\) and \(0.1428\), and the
minimum finite-reference-pool selection ratios are \(0.9876\) and \(0.9911\).
Nevertheless, the sufficient lower bound of \cref{eq:c5-finite-margin} is negative at all 24
transitions in both models. Positive active-space spectra and finite-reference-set ratios therefore do
not establish the continuous relative-score condition. The selected directions remain valid for the
finite numerical experiment, but the transfer proposition cannot be invoked from these quantities alone.
The direct projection estimates below use the selected directions themselves and do not require that
sufficient lower bound to be positive.
The decomposition in \Cref{tab:transfer-margin-example} shows one representative transition.  The Hessian-transfer lower bound is negative even though the corresponding transition is retained; the score-quadrature and active-optimization terms are negligible at this step relative to the nonlinear-to-frozen defect.
\begin{table}[!htbp]
\centering
\scriptsize
\caption{One-step error-budget decomposition at the first calibration transition. The negative margin is a diagnostic: the nonlinear-to-frozen defect dominates the available frozen decrease even though the native step is accepted.}
\label{tab:transfer-margin-example}
\begin{tabularx}{\linewidth}{@{}Xrr@{}}
\toprule
Quantity & Pure \(p=4\), \(8\to9\) & Regularized \(\varepsilon=0.1\), \(8\to9\) \\
\midrule
Frozen decrease & \(1.8403\times10^{-2}\) & \(1.8175\times10^{-2}\) \\
Nonlinear-to-frozen defect & \(6.8347\times10^{-3}\) & \(6.8108\times10^{-3}\) \\
Score quadrature defect & \(9.73\times10^{-13}\) & \(9.95\times10^{-13}\) \\
Active optimization defect & \(4.36\times10^{-7}\) & \(4.48\times10^{-7}\) \\
Total score defect & \(6.8351\times10^{-3}\) & \(6.8112\times10^{-3}\) \\
Selection margin & \(-6.3599\times10^{-2}\) & \(-6.2742\times10^{-2}\) \\
\bottomrule
\end{tabularx}
\end{table}


\paragraph{Fractional-innovation calculation.}
For the benchmark \(\alpha=6\), fixed before evaluating the normalized drops, set
\begin{equation}
 x_N=r_N^2,\qquad
 \gamma_N=\sqrt{x_N-x_{N+1}},\qquad
 c_N=\frac{x_N-x_{N+1}}{x_N^{7/6}},\qquad N=8,\ldots,31.
 \label{eq:numerical-fractional-ratio}
\end{equation}
The minimum over all transitions is \(1.6904\times10^{-2}\) for the pure model and
\(1.0912\times10^{-2}\) for \(\varepsilon=0.1\); both occur at \(25\to26\).
Changing the quadrature order changes these stepwise ratios by less than
\(5.6\times10^{-12}\) relatively. Together with the computed \(q\) and \(\rho\), these
values permit a numerical evaluation of \cref{eq:c5-power,eq:c5-power-energy}.
Writing \(j=N-8\), the resulting finite-range comparison powers are
\[
 \widehat\delta_{8+j}=O(j^{-6}),\qquad
 d_{V,\bullet}(\widehat G_{8+j},u^*)=O(j^{-3}),\qquad
 \|\widehat G_{8+j}-u^*\|_{W^{1,4}}=O(j^{-3/2}),\quad 1\leq j\leq24.
\]
The explicit bound in \cref{eq:c5-power} also includes the starting state \(j=0\).
These powers describe a conditional envelope on the tested finite range, not a limit as \(N\to\infty\).

\begin{table}[!htbp]
\centering
\small
\caption{Fractional-innovation quantities for \(\alpha=6\), computed from every transition \(N\to N+1\), \(N=8,\ldots,31\). The fitted exponent describes the innovation score and is not used to choose the benchmark exponent or its bound.}
\label{tab:fractional-p4}
\begin{tabularx}{\linewidth}{@{}Xrr@{}}
\toprule
Quantity & Pure \(p=4\) & \(\varepsilon=0.1\) \\
\midrule
\(c_{\min}\) over all 24 transitions & \(1.690\times10^{-2}\) & \(1.091\times10^{-2}\) \\
Transition attaining \(c_{\min}\) & \(25\to26\) & \(25\to26\) \\
Fitted slope \(\nu\) of \(\log\gamma_N\) versus \(\log x_N\) & 0.5866 & 0.5898 \\
\(1/(2\nu-1)\) from the fit & 5.77 & 5.57 \\
Fractional bound for \(r_{32}^2\), \(\alpha=6\) & \(3.774\times10^{-2}\) & \(4.110\times10^{-2}\) \\
Fractional bound / computed \(r_{32}^2\) & \(9.510\times10^{3}\) & \(6.821\times10^{3}\) \\
Fractional energy bound at \(N=32\) & \(1.967\times10^{-2}\) & \(2.142\times10^{-2}\) \\
\bottomrule
\end{tabularx}
\end{table}


A least-squares fit of \(\log\gamma_N\) against \(\log x_N\), using all 24 transitions,
gives slopes \(\nu=0.5866\) and \(0.5898\). The relation
\(\alpha=1/(2\nu-1)\) associates these fitted innovation powers with \(5.77\) and
\(5.57\), respectively. This fit concerns \(\gamma_N=\theta_NS_N/\ell_N\), not \(S_N\)
alone; it does not independently establish either factor in \cref{eq:c5-power-factors}.

The worst-step constants still yield conservative terminal projection bounds:
\(3.77\times10^{-2}\) and \(4.11\times10^{-2}\), approximately
\(9.5\times10^3\) and \(6.8\times10^3\) times the computed projection errors.
The fractional bound is slightly smaller than the source bound for the pure model and slightly
larger for the regularized model. Positive drops on a finite sequence also give a positive
minimum normalized ratio for any fixed exponent. The computation establishes the stated numerical
envelope and its constants, but does not uniquely identify sixth-order decay or verify an
asymptotic mechanism. Such a conclusion would require bounds uniform over increasing index ranges.

\paragraph{Continuation beyond the calibration range.}
To test whether the calibrated constants remain informative immediately beyond the range on which
they were fixed, we continued the two trajectories from \(N=32\) to \(N=64\), keeping the same
candidate pools, previously computed iterates, and source coefficients and drawing no further random atoms. The
comparison exponent and all constants were kept at their \(N=8\)--\(32\) values. Most transitions in
this subsequent finite range satisfy the normalized-innovation sufficient condition with those constants.
At quadrature order \(20\), it holds for 27 of the 32 continuation transitions in the pure model
\(27/32=84.4\%\) and 28 of 32 in the regularized model
\(28/32=87.5\%\); only five and four transitions, respectively, fall
below the prescribed threshold. These exceptions concern the sufficient condition, not the validity
of the theoretical implication. In both models, the first such transition is \(N=33\to34\), and the
minimum ratio \(c_N/c_0\) over the continuation is \(0.156\) and \(0.187\), respectively; the
terminal squared projection errors are \(4.770\times10^{-7}\) and \(4.626\times10^{-7}\). The
order-16 cross-check changes these reported values only in the displayed last digits. The
continuation therefore shows that the calibrated condition remains satisfied for
most transitions, together with a substantial error decrease and agreement with the algebraic
identities on the subsequent states. The few exceptions prevent applying the same fixed-constant
theorem over the entire continuation range. The stated finite-range envelope therefore retains
its calibration range $N=8$--32, while the continuation provides supporting stepwise numerical
evidence without establishing an asymptotic rate.

\begin{table}[!htbp]
\centering\small
\caption{Continuation from $N=32$ to $N=64$ (32 transitions) with constants fixed on $N=8$--$32$. Pass counts refer to the normalized-innovation sufficient condition; the first failed transition starts at the reported $N$.}
\label{tab:continuation-window}
\begin{tabular}{@{}lrrrrr@{}}
\toprule
Model & New steps & Pass & First failed $N$ & $\min c_N/c_0$ & $r_{64}^2$ \\
\midrule
Pure $p=4$ & 32 & 27 & 33 & 0.156 & $4.770\times10^{-7}$ \\
$\varepsilon=0.1$ & 32 & 28 & 33 & 0.187 & $4.626\times10^{-7}$ \\
\bottomrule
\end{tabular}
\end{table}

\FloatBarrier

